\section{Euler logarithms and prime traces}
\label{v4-arith:ks:logL-not-cyclic}

An Euler product already has a concrete algebraic carrier: the free
commutative monoid on the primes. Its logarithm singles out prime powers,
and an energy derivation supplies the factors $\log p$. These two
operations must be distinguished before seeking a Hall-algebra realization.

Let $\mathscr D_K$ be the algebra of formal Dirichlet series over a
commutative $\mathbb Q$-algebra $K$. An element is a family
$a=\sum_{n\geq1}a_n[n]$, with multiplication $[m][n]=[mn]$;
the coefficient of $[n]$ in $ab$ is $\sum_{d\mid n}a_db_{n/d}$.
We use the coefficient topology, or equivalently the filtration by the
ideals of series supported on $n\geq N$. If $a_1=0$, then the coefficient
of $[n]$ in $a^j$ vanishes for $2^j>n$, so formal logarithms and
exponentials converge in this topology.

\begin{proposition}[Euler logarithm and energy derivation]
\label{v4-arith:obs:logL-not-cyclic}
\ClaimStatusProvedHere{}
In $\mathscr D_{\mathbb Q}$ one has
\begin{equation}
 Z=\prod_p(1-[p])^{-1}=\sum_{n\geq1}[n],
 \qquad
 \log Z=\sum_p\sum_{k\geq1}\frac{[p^k]}{k}.
 \label{v4-arith:ks:euler-log}
\end{equation}
After extending coefficients to $\mathbb R$, the continuous derivation
$D[n]=(\log n)[n]$ satisfies
\begin{equation}
 D\log Z=\sum_{n\geq1}\Lambda(n)[n],\qquad
 \Lambda(n)=
 \begin{cases}\log p&n=p^k,\ k\geq1,\\0&\text{otherwise.}\end{cases}
 \label{v4-arith:ks:euler-derivation}
\end{equation}
For $\operatorname{Re}s>1$, their convergent analytic realizations are
\begin{equation}
 \log\zeta(s)=\sum_{p,k\geq1}\frac{p^{-ks}}{k},
 \qquad
 -\frac{\zeta'(s)}{\zeta(s)}
   =\sum_{p,k\geq1}(\log p)p^{-ks}.
 \label{v4-arith:ks:euler-analytic}
\end{equation}
The logarithm is the holomorphic branch tending to zero as
$\operatorname{Re}s\longrightarrow+\infty$.
\end{proposition}

\begin{proof}
Unique prime factorization gives exactly one contribution to every
coefficient of the product in \eqref{v4-arith:ks:euler-log}.
Taking its logarithm coefficient by coefficient gives
$-\log(1-[p])=\sum_{k\geq1}[p^k]/k$.
The identity $\log(mn)=\log m+\log n$ proves the Leibniz rule for $D$.
Since $D[p^k]=k\log p\,[p^k]$, applying $D$ cancels the denominator $k$.

On a compact subset of $\operatorname{Re}s>1$, choose $\sigma_0>1$
below all real parts. Then
\[
 \sum_{p,k\geq1}\frac{|p^{-ks}|}{k}
 \leq\frac{1}{1-2^{-\sigma_0}}\sum_{n\geq2}n^{-\sigma_0}<\infty.
\]
The differentiated series has the same bound with
$\sum_{n\geq2}(\log n)n^{-\sigma_0}$ in place of the final sum.
Thus both series converge locally uniformly. Exponentiation identifies
the first series with the Euler product for $\zeta$ and shows that this
product is nonzero. Differentiation proves the second identity.
\end{proof}

The coefficient of $n^{-s}$ in $\log\zeta$ is therefore $1/k$ for
$n=p^k$, and zero otherwise. It is rational. The factor $\log p$
belongs to the differentiated series. Neither assertion involves a
cyclic-cohomology construction.

\begin{proposition}[local determinant factors]
\label{v4-arith:ks:local-determinant-traces}
\ClaimStatusProvedHere{}
For each prime $p$, let $T_p$ be a complex endomorphism of a Hermitian
space of dimension at most $d$, where $d$ is fixed. Suppose
$\|T_p\|\leq p^\theta$ for a fixed $\theta\geq0$. On
$\operatorname{Re}s>1+\theta$ the product
\[
 L(s)=\prod_p\det(1-p^{-s}T_p)^{-1}
\]
is holomorphic and nonzero, and
\begin{align}
 \log L(s)&=\sum_{p,k\geq1}
      \frac{\operatorname{Tr}(T_p^k)}{k}p^{-ks},
      \label{v4-arith:ks:matrix-euler-log}\\
 -\frac{L'(s)}{L(s)}&=\sum_{p,k\geq1}
      (\log p)\operatorname{Tr}(T_p^k)p^{-ks}.
      \label{v4-arith:ks:matrix-euler-derivative}
\end{align}
Here $\operatorname{Tr}$ is the ordinary, unnormalized matrix trace,
and the logarithm tends to zero as the real part of $s$ tends to infinity.
\end{proposition}

\begin{proof}
For $|q|\|T\|<1$, Jacobi's determinant formula gives
\[
 \frac{d}{dq}\bigl[-\log\det(1-qT)\bigr]
   =\operatorname{Tr}\bigl(T(1-qT)^{-1}\bigr)
   =\sum_{k\geq1}\operatorname{Tr}(T^k)q^{k-1}.
\]
Integration from $q=0$ fixes the constant and gives the local logarithm.
The estimate $|\operatorname{Tr}(T_p^k)|\leq d p^{k\theta}$ reduces
normal convergence, including differentiation, to the preceding proof
with $\operatorname{Re}s-\theta$ in place of $\operatorname{Re}s$.
\end{proof}

These formulas give an arithmetic use for a proposed graded carrier:
its trace must produce the specified local factors. For $T_p=1$, that
carrier can be constructed without a quiver.

\begin{proposition}[the prime monoid and its thermal trace]
\label{v4-arith:ks:prime-fock-trace}
\ClaimStatusProvedHere{}
Let $\mathcal P=\mathbb Q[x_p:p\text{ prime}]$, with each monomial
involving finitely many variables. Declare its monomials orthonormal
after extension to $\mathbb C$, and complete to a Hilbert space
$\mathcal F$. The assignment
\[
 x^\nu\longmapsto\varepsilon_{\prod_p p^{\nu_p}}
\]
is a unitary identification $\mathcal F\cong\ell^2(\mathbb N_{\geq1})$.
The diagonal energy operator
\[
 Hx^\nu=\left(\sum_p\nu_p\log p\right)x^\nu
\]
is nonnegative and self-adjoint on
$\{\sum_n a_n\varepsilon_n:\sum_n(\log n)^2|a_n|^2<\infty\}$.
For $\operatorname{Re}s>1$,
\begin{equation}
 \operatorname{Tr}(e^{-sH})=\zeta(s),\qquad
 \operatorname{Tr}(He^{-sH})=-\zeta'(s).
 \label{v4-arith:ks:prime-thermal-trace}
\end{equation}
For real $\beta>1$, the normalized state
$\langle A\rangle_\beta=\operatorname{Tr}(Ae^{-\beta H})/\zeta(\beta)$,
whenever the numerator is defined, satisfies
\[
 \langle H\rangle_\beta=-\frac{\zeta'(\beta)}{\zeta(\beta)},
 \qquad
 \langle H^2\rangle_\beta-\langle H\rangle_\beta^2
   =\sum_{p,k\geq1}k(\log p)^2p^{-k\beta}.
\]
\end{proposition}

\begin{proof}
The monomial basis is indexed bijectively by positive integers through
unique factorization. Under this identification, $H\varepsilon_n
=(\log n)\varepsilon_n$. A real diagonal multiplication operator on
its maximal square-summability domain is self-adjoint; this follows
directly by testing the adjoint against each basis vector.
The singular values of $e^{-sH}$ are $n^{-\operatorname{Re}s}$, so it
is trace class on the stated half-plane. The same argument applies to
$H^je^{-sH}$ for every fixed $j\geq0$. Its diagonal trace proves
\eqref{v4-arith:ks:prime-thermal-trace} and permits two differentiations.
The quotient rule then identifies the variance with
$(\log\zeta)''(\beta)$, and \eqref{v4-arith:ks:euler-analytic} evaluates it.
\end{proof}

For a finite set of primes $F$, the same construction on
$\mathcal P_F=\mathbb Q[x_p:p\in F]$ has partition function
$\zeta_F(s)=\prod_{p\in F}(1-p^{-s})^{-1}$ for
$\operatorname{Re}s>0$. Multiplication by $x_p$ extends to the isometry
$\varepsilon_n\mapsto\varepsilon_{pn}$. Thus the monoid action gives
the prime shifts in the Bost--Connes representation. The displayed
trace is a Hilbert-space trace; a map to the arithmetic Hochschild
trace class $\mathcal V^{\mathrm{prim}}$ is additional data.

\section{Hall multiplication and the BPS filtration}
\label{v4-arith:ks:CoHA-conjectural}

A Hall product counts extensions of representations. A monoid product
multiplies their prime labels. Even for one vertex these operations
produce different algebras, and the distinction persists when their
graded characters resemble Euler products.

\begin{proposition}[the zero-loop and one-loop calculations]
\label{v4-arith:ks:one-vertex-coha}
\ClaimStatusProvedHere{}
Let $Q_m$ have one vertex and $m\in\{0,1\}$ loops, with zero potential.
Its rational CoHA has dimension-$n$ summand
$\mathcal H^{(m)}_n=\mathbb Q[x_1,\ldots,x_n]^{\mathfrak S_n}$.
A homogeneous polynomial of polynomial degree $j$ has cohomological
degree $2j+(1-m)n^2$. Write $e_k=x^k\in\mathcal H^{(m)}_1$.
Then
\[
 \mathcal H^{(0)}\cong\bigwedge_{\mathbb Q}(e_0,e_1,e_2,\ldots),
 \quad |e_k|=2k+1,
 \qquad
 \mathcal H^{(1)}\cong\mathbb Q[e_0,e_1,e_2,\ldots],
 \quad |e_k|=2k.
\]
Every generator has dimension degree one.
\end{proposition}

\begin{proof}
The representation space is affine and equivariantly contractible,
so its equivariant cohomology is $H^*(B\mathrm{GL}_n,\mathbb Q)$.
The Hall correspondence chooses an invariant subspace. Localization
at the diagonal torus gives the shuffle formula
\begin{equation}
 f*g=\sum_{\substack{I\sqcup J=\{1,\ldots,a+b\}\\|I|=a}}
 f(x_I)g(x_J)\prod_{i\in I,\,j\in J}(x_j-x_i)^{m-1}.
 \label{v4-arith:ks:one-vertex-shuffle}
\end{equation}
The numerator is the Euler class of the $m$ arrow blocks and the
denominator the Euler class of the tangent space to the Grassmannian.
This is the specialization of Kontsevich--Soibelman,
\emph{Cohomological Hall algebra, exponential Hodge structures and
motivic Donaldson--Thomas invariants}, Theorem~2, pp.~11--12.

For $m=0$, the product of $n$ dimension-one elements is their
alternant divided by $\prod_{i<j}(x_j-x_i)$. Repeated exponents give
zero. For $0\leq k_1<\cdots<k_n$, it is the Schur polynomial of
partition $(k_n-n+1,k_{n-1}-n+2,\ldots,k_1)$. These form a basis of
the symmetric polynomials. Hence the exterior-algebra map is an
isomorphism. In particular,
\[
 e_0*e_0=0,\qquad
 e_0*e_1=\frac{x_2-x_1}{x_2-x_1}=1\in\mathcal H^{(0)}_2,
 \qquad e_1*e_0=-1\in\mathcal H^{(0)}_2.
\]
The constants on the right lie in dimension degree two and are not
the unit, which lies in dimension degree zero.

For $m=1$, the shuffle factor is one. Products of the $e_k$ give the
monomial symmetric basis, multiplied by the factorials of repeated
exponents. Those factorials are invertible over $\mathbb Q$, proving
the polynomial-algebra assertion. This argument also explains its
coefficient restriction: over $\mathbb Z$, for example,
$e_0*e_0=2\cdot1\in\mathcal H^{(1)}_2$.
\end{proof}

An edgeless quiver with a vertex for each prime in a nonempty finite
set $F$ therefore has odd square-zero classes already in each
one-vertex subalgebra. Its CoHA cannot be $\mathcal P_F$.
For $F=\varnothing$, there is only the zero representation, so both
the CoHA and $\mathcal P_{\varnothing}$ are $\mathbb Q$.
For the disjoint union of Jordan quivers, the rational Hall algebra
is polynomial on generators $e_{p,k}$ for $p\in F$, $k\geq0$.
There is, however, a precise prime-monoid comparison:
\begin{equation}
 \mathcal P_F\ \xrightarrow{\ x_p\mapsto e_{p,0}\ }\
 \mathcal H_F^{(1)},\qquad
 \mathcal H_F^{(1)}/(e_{p,k}:k>0)\cong\mathcal P_F.
 \label{v4-arith:ks:jordan-prime-maps}
\end{equation}
The first map identifies the cohomological-degree-zero subalgebra;
the second kills the positive-degree generators. The maps are
compatible with inclusions of finite sets of primes. The full local
graded character, as a series in $\mathbb Q[[z,q]]$ with $q$ recording
half the cohomological degree, is
\[
 \sum_{n,j\geq0}\dim\bigl(\mathcal H^{(1)}_{n,2j}\bigr)z^nq^j
   =\prod_{k\geq0}(1-zq^k)^{-1}.
\]
Its coefficient of $q^0$ is $(1-z)^{-1}$, the prime Euler factor
after $z=p^{-s}$. Forgetting the $q$-grading instead gives infinite
multiplicities even in dimension degree one. Thus
\eqref{v4-arith:ks:jordan-prime-maps} specifies the sector on which
Proposition~\ref{v4-arith:ks:prime-fock-trace} applies.

To include a potential, fix a finite quiver $Q$ with symmetric arrow
matrix and a polynomial cyclic potential
$W\in\mathbb C Q/[\mathbb C Q,\mathbb C Q]_{\mathrm{vect}}$.
For a dimension vector $\mathbf d\in\mathbb N^{Q_0}$, put
\[
 X_{\mathbf d}=\bigoplus_{a:i\to j}
     \operatorname{Hom}(\mathbb C^{d_i},\mathbb C^{d_j}),\qquad
 G_{\mathbf d}=\prod_i\mathrm{GL}_{d_i}(\mathbb C),\qquad
 \mathfrak M_{\mathbf d}=[X_{\mathbf d}/G_{\mathbf d}].
\]
The Euler form is
$\chi(\mathbf d,\mathbf e)=\sum_i d_ie_i-\sum_{a:i\to j}d_ie_j$,
so $\dim\mathfrak M_{\mathbf d}=-\chi(\mathbf d,\mathbf d)$.
Evaluation followed by matrix trace defines the invariant function
$f_{\mathbf d}=\operatorname{Tr}(W)$ on this stack and on its affine
coarse moduli space $M_{\mathbf d}=X_{\mathbf d}/\!/G_{\mathbf d}$.
Assume $\operatorname{crit}(f_{\mathbf d})\subset f_{\mathbf d}^{-1}(0)$
for every $\mathbf d$; zero and homogeneous positive-degree potentials
satisfy this condition.

We use rational monodromic mixed Hodge structures and the monodromic
vanishing-cycle functor $\phi_f^{\mathrm{mon}}$, normalized to be
perverse exact and to satisfy $\phi_0^{\mathrm{mon}}=\mathrm{id}$.
Let $\mathbb L=H_c^*(\mathbb A^1,\mathbb Q)$, in cohomological
degree two. Its monodromic square root is
$\mathbb L^{1/2}=H_c^*(\mathbb A^1,\phi_{x^2}^{\mathrm{mon}}
\mathbb Q_{\mathbb A^1})$, in degree one, with the square-root
isomorphism fixed by the complex orientation. On a smooth space or smooth quotient
stack $Y$, write
$\mathrm{IC}_Y^{\mathrm{vir}}=\mathbb L^{-\dim Y/2}\otimes\mathbb Q_Y$.
On a singular irreducible coarse moduli space, use the intersection
complex with the same normalization on its smooth locus and
intermediate extension. These conventions include the shifts in the
one-vertex calculation.

\begin{definition}[critical cohomological Hall algebra]
\label{v4-arith:thm:KS-CoHA}
The underlying dimension-graded, cohomologically graded object is
\begin{equation}
 \mathcal H_{Q,W}
 =\bigoplus_{\mathbf d}
 \left[
 H_c^*(\mathfrak M_{\mathbf d},
       \phi_{f_{\mathbf d}}^{\mathrm{mon}}
       \mathrm{IC}_{\mathfrak M_{\mathbf d}}^{\mathrm{vir}})
 \right]^\vee.
 \label{v4-arith:ks:normalized-coha}
\end{equation}
Here $\vee$ denotes graded duality, including duality of Hodge
structures, and equivariant cohomology is defined by finite-dimensional
approximations to the classifying spaces. Its associative product is
the vanishing-cycle pullback and proper pushforward along the stack
of short exact sequences, with the virtual shifts displayed above.
The equality of trace potentials on a block upper-triangular
representation and on its subquotients supplies the required
Thom--Sebastiani map. The zero representation supplies the unit.
This is the critical CoHA of Kontsevich--Soibelman, \S7, in the
normalization of Davison--Meinhardt, equation~(1) and \S5.2.
\end{definition}

Choose a stability condition with all central charges $\zeta_i=\sqrt{-1}$.
Every representation is semistable of slope zero, since the slope
of $\mathbf d\ne0$ is
$-\operatorname{Re}(\sum_i\zeta_i d_i)/
\operatorname{Im}(\sum_i\zeta_i d_i)=0$.
This stability is slope-generic: the antisymmetrization of $\chi$
vanishes on all dimension vectors because $Q$ is symmetric.
Thus the single-slope integrality theorem applies to the full algebra
\eqref{v4-arith:ks:normalized-coha}.

For $\mathbf d\ne0$, define a BPS sheaf on $M_{\mathbf d}$ by
\[
 \mathcal B_{\mathbf d}=
 \begin{cases}
 \phi_{f_{\mathbf d}}^{\mathrm{mon}}
       \mathrm{IC}_{M_{\mathbf d}}^{\mathrm{vir}}
       &\text{if the stable locus is nonempty},\\
 0&\text{otherwise},
 \end{cases}
 \qquad
 B=\bigoplus_{\mathbf d\ne0}H^*(M_{\mathbf d},\mathcal B_{\mathbf d}).
\]
At this stability, stable representations are precisely the simple
ones. Finally put
\[
 U=H^*(B\mathbb C^*,\mathbb Q)_{\mathrm{vir}}\otimes B
   =\mathbb L^{1/2}\otimes\mathbb Q[u]\otimes B,
 \qquad |u|=2.
\]
Here $u$ is the first Chern class of the universal line bundle and
has dimension degree zero. In particular, the BPS
cohomology $B$, the shifted BPS Lie algebra $\mathbb L^{1/2}\otimes B$,
and the full generator object $U$ are distinct.

\begin{theorem}[Davison--Meinhardt integrality and PBW]
\label{v4-arith:thm:davison-integrality}
\ClaimStatusProvedElsewhere{}
For the preceding data define, modulo two,
\[
 \tau(\mathbf d,\mathbf e)=\chi(\mathbf d,\mathbf e)
       +\chi(\mathbf d,\mathbf d)\chi(\mathbf e,\mathbf e).
\]
Choose a bilinear form $\psi$ with
$\psi(\mathbf d,\mathbf e)+\psi(\mathbf e,\mathbf d)
=\tau(\mathbf d,\mathbf e)$, and twist the Hall product by
$a*_\psi b=(-1)^{\psi(\mathbf d,\mathbf e)}a*b$.
Write $\mathcal H^\psi$ for the resulting algebra. It has an
exhaustive multiplicative ambient perverse filtration $P$ and a
canonical generator embedding $\iota:U\hookrightarrow\mathcal H^\psi$
such that the symmetrized Hall-product map
\begin{equation}
 \operatorname{PBW}:\operatorname{Sym}_{\mathrm{sup}}(U)
       \xrightarrow{\ \sim\ }\mathcal H^\psi
 \label{v4-arith:ks:pbw-object-map}
\end{equation}
is an isomorphism of dimension-graded, cohomologically graded
monodromic mixed Hodge structures. The associated graded
$\operatorname{gr}_P\mathcal H^\psi$ is a free supercommutative
algebra. The positive-dimension part of $P_1\mathcal H^\psi$ is
$\mathbb L^{1/2}\otimes B$, and is closed under the Hall
supercommutator.
\end{theorem}

Here $\operatorname{Sym}_{\mathrm{sup}}$ is polynomial on the even
part and exterior on the odd part, with the cohomological Koszul sign.
The form $\tau$ is alternating modulo two, so a form $\psi$ exists
by choosing an ordering of a basis of $(\mathbb Z/2)^{Q_0}$.
Bilinearity makes the twisted product associative. The theorem is
Davison--Meinhardt, \emph{Cohomological Donaldson--Thomas theory of
a quiver with potential and quantum enveloping algebras},
Theorems~A and~C, pp.~5--6 and 10--11; the Lie subalgebra is
\S1.7, equation~(22).

For clarity, $P$ uses the semisimplification map
$p_{\mathbf d}:\mathfrak M_{\mathbf d}\to M_{\mathbf d}$.
One approximates this map by the proper maps from moduli spaces
of sufficiently framed representations to $M_{\mathbf d}$, with
their virtual shifts. Perverse truncation of these direct images
on the ambient space $M_{\mathbf d}$, followed by cohomology and
stabilization in each degree, defines $P$; equivalently one uses
the dual compact-support construction with truncation
$\tau_{\geq-p}$ in Davison--Meinhardt, \S5.3, pp.~48--49.
The summand $\mathbb L^{1/2}u^k\otimes B$ supplies generators of
perverse degree $2k+1$. Supercommutativity of the associated graded
implies $[P_a,P_b]\subset P_{a+b-1}$, and hence the assertion about
$P_1$. No integral lattice or coefficient-change theorem is asserted
by this rational Hodge-theoretic construction.

\begin{proposition}[the multiplicative PBW criterion]
\label{v4-arith:ks:pbw-multiplicative-criterion}
\ClaimStatusProvedHere{}
The map \eqref{v4-arith:ks:pbw-object-map} is an algebra isomorphism
if and only if the embedded generators $\iota(U)$ pairwise
supercommute in $\mathcal H^\psi$. An ordinary polynomial algebra
on these homogeneous generators additionally requires $U$ to have
no odd part.
\end{proposition}

\begin{proof}
If the generators supercommute, the universal property gives an
algebra map from $\operatorname{Sym}_{\mathrm{sup}}(U)$ to
$\mathcal H^\psi$. On a product of $n$ generators it equals the
signed average of the $n!$ ordered Hall products, hence equals the
PBW map and is bijective. Conversely, an algebra map extending
$\iota$ preserves the defining supercommutation relations.
Over $\mathbb Q$ a nonzero odd generator in the symmetric
superalgebra has square zero, whereas an ordinary polynomial
algebra over $\mathbb Q$ is reduced.
\end{proof}

The commutator condition can fail even for a symmetric quiver with
a homogeneous cubic potential.

\begin{example}[a noncommutative symmetric critical CoHA]
\label{v4-arith:ks:tripled-a2}
Start with $Q=(1\xrightarrow{a}2)$, add the reverse arrow $a^*$
and loops $\omega_1,\omega_2$, and call the resulting symmetric
quiver $\widetilde Q$. With path multiplication given by composition
from right to left, take the cyclic potential
\[
 \widetilde W=(aa^*-a^*a)(\omega_1+\omega_2).
\]
In dimension $(1,1)$ its trace is
$(\omega_2-\omega_1)aa^*$. For the sign twist choose the original
Euler form
$\psi(\mathbf d,\mathbf e)
=d_1e_1+d_2e_2-d_1e_2\pmod2$.
It obeys the required relation because
$\chi_{\widetilde Q}(\mathbf d,\mathbf e)
=-d_1e_2-d_2e_1$.

Dimensional reduction identifies the resulting critical CoHA with
the preprojective CoHA of $Q$, with this multiplication and twist.
The degree-zero BPS Lie algebra is
$\mathfrak n^-_{\mathfrak{sl}_3}$: this is Davison,
\emph{BPS Lie algebras and the less perverse filtration on the
preprojective CoHA}, \S3.2, equation~(35), and Theorem~6.6,
pp.~33--34. That theorem concerns the Lie algebra embedded in the
actual Hall algebra, not just an associated graded algebra.
Under its simple-root normalization, the two vertex classes
$\alpha_1,\alpha_2$ correspond to $E_{21},E_{32}$. Consequently
\[
 [\alpha_2,\alpha_1]\ne0,
 \qquad [E_{32},E_{21}]=E_{31}\ne0.
\]
Both classes have cohomological degree zero, so this is an ordinary
commutator of even elements. The critical CoHA is not
supercommutative, although Theorem~\ref{v4-arith:thm:davison-integrality}
gives its PBW decomposition. Thus that theorem cannot imply a
polynomial-algebra assertion for an arbitrary arithmetic potential.
\end{example}

\section{Cyclic potentials and the arithmetic comparison}
\label{v4-arith:ks:arithmetic-potential-comparison}

The logarithmic Euler series does have a cyclic realization, but its
critical locus tests whether it can serve as a Hall potential.
For an associative algebra $A$, degree-zero cyclic \emph{homology}
is $\mathrm{HC}_0(A)=A/[A,A]_{\mathrm{vect}}$.
Over a field, degree-zero cyclic cohomology consists of linear
traces $\lambda:A\to K$ satisfying $\lambda(ab)=\lambda(ba)$.
These constructions neither adjoin coefficients nor supply a map
from an arbitrary Dirichlet series into a path algebra.

\begin{proposition}[the completed logarithmic potential]
\label{v4-arith:ks:logarithmic-potential-obstruction}
\ClaimStatusProvedHere{}
For a finite set of primes $F$, let $Q_F$ be the disjoint union of
one-loop quivers, with loop $x_p$ at vertex $p$. Complete its complex
path algebra in the arrow ideal. In the quotient by the closure of
the commutator subspace, the cyclic series
\[
 W_F=\sum_{p\in F}\sum_{k\geq1}\frac{x_p^k}{k}
\]
is well defined. On tuples of finite matrices with $\|X_p\|<1$,
\[
 \operatorname{Tr}(W_F)(X)
   =-\sum_{p\in F}\log\det(1-X_p).
\]
In particular, at dimension one at each vertex and
$X_p=p^{-s}$ it equals $\log\zeta_F(s)$ for
$\operatorname{Re}s>0$. Its completed Jacobi algebra is zero.
\end{proposition}

\begin{proof}
There are no paths between different vertices, so the completed
path algebra is $\prod_{p\in F}\mathbb C[[x_p]]$.
Convergence and the trace identity follow from the local determinant
calculation. The cyclic derivative removes one occurrence of the
specified arrow and sums over its positions. Hence
\[
 \frac{\partial W_F}{\partial x_p}
   =\sum_{k\geq1}x_p^{k-1}=(1-x_p)^{-1}.
\]
This element is a unit in the corner with identity $e_p$.
The closed two-sided ideal generated by the cyclic derivatives
therefore contains every $e_p$ and hence the identity $\sum_{p\in F}e_p$.
Its quotient, the completed Jacobi algebra, is zero.
Equivalently, in every nonzero dimension vector the trace function
has no critical point on the matrix domain: its derivative in a variation $\delta X_p$ is
$\operatorname{Tr}((1-X_p)^{-1}\delta X_p)$, a nonzero covector.
\end{proof}

This series lies in a completed cyclic path algebra, whereas
Theorem~\ref{v4-arith:thm:davison-integrality} is stated for a
polynomial potential. The obstruction is already visible for every
polynomial truncation containing its linear term: in the formal
neighborhood of the zero representation, its cyclic derivative is
a unit and its vanishing cycles are zero in positive dimension.
Thus the displayed logarithmic realization supplies no nonzero
local BPS object at the origin.

Removing the low-degree terms changes the critical problem in a
computable way. At one vertex, for $m\geq2$, put
\[
 W^{\geq m}=\sum_{k\geq m}\frac{x^k}{k},\qquad
 \frac{\partial W^{\geq m}}{\partial x}
   =\frac{x^{m-1}}{1-x}.
\]
Its completed Jacobi algebra is
$\mathbb C[[x]]/(x^{m-1})$: it is $\mathbb C$ for $m=2$ and
the dual numbers for $m=3$. The scalar trace at $x=p^{-s}$ is now
$-\log(1-p^{-s})-\sum_{1\leq k<m}p^{-ks}/k$.
The missing terms must therefore be supplied by a separate
construction if this singularity is used in an Euler-factor model.

\begin{proposition}[arithmetic comparison criterion]
\label{v4-arith:thm:arithmetic-CoHA-conjectural}
\ClaimStatusConditional{}
Suppose that for every finite set of primes $F$ one has a finite
symmetric complex quiver with polynomial potential $(Q_F,W_F)$
satisfying the hypotheses above, and a dimension-graded rational
object $V_F$ with a specified isomorphism
\[
 \eta_F:V_F\xrightarrow{\ \sim\ }B_{Q_F,W_F}
\]
of cohomologically graded monodromic mixed Hodge structures.
Then PBW gives an object isomorphism
\[
 \operatorname{Sym}_{\mathrm{sup}}
    (\mathbb L^{1/2}\otimes\mathbb Q[u]\otimes V_F)
       \xrightarrow{\ \sim\ }\mathcal H_{Q_F,W_F}^{\psi}.
\]
It is an algebra isomorphism precisely when the corresponding
embedded generator tower supercommutes. If these generators
supercommute and are all even, it is an ordinary polynomial-algebra
isomorphism.
\end{proposition}

\begin{proof}
Apply Theorem~\ref{v4-arith:thm:davison-integrality} through
$\mathbb L^{1/2}\otimes\mathbb Q[u]\otimes\eta_F$ and then apply
Proposition~\ref{v4-arith:ks:pbw-multiplicative-criterion}.
\end{proof}

An arithmetic identification with $\mathcal V^{\mathrm{prim}}$
requires more than the objects $V_F$. It requires a realization of
that Hochschild trace class in the same graded coefficient category,
a comparison with the $V_F$, and maps compatible with enlarging $F$.
To obtain an Euler product, one must also specify an energy or
graded-character operation and prove that it intertwines the
comparison with the local traces of
Proposition~\ref{v4-arith:ks:local-determinant-traces}.
The polynomial algebra in that comparison contains the entire
$\mathbb Q[u]$ tower. Restricting to a single generator per prime
needs a specified subalgebra or quotient, as in
\eqref{v4-arith:ks:jordan-prime-maps}.

The construction of such quiver data and compatible comparison maps
for the arithmetic Hochschild trace remains open. The prime Fock
trace proves the Euler product on an explicit carrier; the
zero-loop calculation excludes an edgeless CoHA on a nonempty prime
set as that carrier;
the logarithmic-potential calculation excludes its most direct
completed critical realization; and the tripled-$A_2$ calculation
shows why a PBW theorem alone cannot remove the multiplication
obstruction.
